% !TeX TS-program = xelatex
% 虚构演示简历 — 仅用于 hellojobs 测试数据
\documentclass{resume-photo}
\usepackage{xeCJK}
\setCJKmainfont{Noto Serif CJK SC}
\usepackage{fontspec}
\usepackage{enumitem}
\setlist[itemize]{itemsep=0pt, topsep=1pt, parsep=0pt, partopsep=0pt}

\ResumeName{吴昊}

\begin{document}

\ResumeContacts{
  138-0008-0008,%
  \ResumeUrl{mailto:wuhao@uestc.edu.cn}{wuhao@uestc.edu.cn},%
  \textnormal{电子科技大学 | 计算机科学与技术 · 学士 | 2021—2025}%
}

\ResumeTitle

\noindent 游戏客户端与图形，熟悉 C++17、Unreal 蓝图与性能分析；参与手游战斗同步与延迟补偿模块开发。注重可观测性与文档沉淀，习惯在评审阶段拆解风险；参与线上复盘与跨团队联调，英语 CET-6，可阅读官方文档与 RFC（演示数据）。

\section{教育背景}
\ResumeItem
{电子科技大学}
[\textnormal{计算机科学与工程学院 | 计算机科学与技术 | 学士}]
[2021.09—2025.06]

\begin{itemize}
  \item 图形学、游戏引擎选修为主。
\end{itemize}


\ResumeItem
{实验室与助教}
[\textnormal{本科生科研 / 课程助教}]
[2022—2025]

\begin{itemize}
  \item 进入校级重点实验室参与课题，负责数据采集、清洗与可视化看板搭建。
  \item 担任《数据结构》《操作系统》课程助教，批改作业 200+ 份，答疑课时 40h+。
  \item 组织期中复习串讲与上机辅导，课程评教得分位于前 10\%（演示）。
  \item 整理课程 FAQ 与实验踩坑文档，被后续年级沿用。
\end{itemize}



\section{专业技能}
\begin{itemize}
  \item 熟悉 C++、Unreal Engine 5、内存与渲染性能剖析。
  \item 了解 UDP 状态同步、插值与客户端预测。
  \item 掌握 Git Perforce 协作与 CrashSight 符号化。
  \item 熟悉 Linux 日常运维与 Shell，具备日志检索与 CPU/内存突增初步定位能力。
  \item 掌握 Git / CI 与 Code Review，了解 Docker 与 Kubernetes 基本编排与滚动发布。
  \item 了解 OAuth2 / JWT、接口幂等与 REST 规范，具备单元测试与集成测试习惯（JUnit/pytest 等）。
  \item 能阅读英文技术文档，具备跨团队沟通与需求澄清能力（测试简历）。
\end{itemize}

\section{实习经历}
\ResumeItem
{某头部手游工作室}
[\textnormal{客户端开发实习生}]
[2024.06—2024.12]

\begin{itemize}
  \item \textbf{职责}: PVP 技能表现与网络同步问题排查。
  \item \textbf{成果}: 高延迟场景下技能判定投诉量下降约 40\%。
\end{itemize}


\ResumeItem
{百度 | 移动生态}
[\textnormal{后端研发实习生}]
[2023.06—2023.11]

\begin{itemize}
  \item \textbf{检索}: 参与搜索相关性特征工程与离线任务调度稳定性治理。
  \item \textbf{存储}: 优化 HBase Get 批量模式，P95 延迟下降 19\%。
  \item \textbf{平台}: 熟悉内部 BVC 构建与发布流水线。
  \item \textbf{算法对接}: 与算法同学联调 A/B 实验埋点与分流一致性校验。
  \item \textbf{文档}: 维护接口字典与错误码说明，降低联调沟通成本。
\end{itemize}



\section{项目经历}
\ResumeItem
{Mini 物理引擎实验}
[\textnormal{UE5 + C++}]
[2023.10—2024.04]

\begin{itemize}
  \item 实现简化刚体碰撞与固定步长积分，用于教学 Demo。
\end{itemize}


\ResumeItem
{多租户 SaaS 计费引擎}
[\textnormal{订阅制 + 用量计费（演示）}]
[2023.09—2024.02]

\begin{itemize}
  \item \textbf{模型}: 设计套餐、用量累计、账单周期与 proration 规则。
  \item \textbf{一致}: 使用 Outbox + 消息表保证计费事件与财务对账一致。
  \item \textbf{性能}: 批价任务分片并行，账单生成时间缩短 48\%。
  \item \textbf{测试}: 构造边界用例库 150+，覆盖闰月与升级降级场景。
\end{itemize}



\section{个人荣誉}
\begin{itemize}
  \item 腾讯游戏高校创意制作大赛 赛区优秀奖
  \item 电子科大 优秀共青团员
  \item 全国计算机等级考试四级（网络工程师）（演示）
  \item 校级「优秀学生干部 / 优秀团员」或技术社团骨干（综合测评前 15\%）
  \item 开源贡献：向知名项目提交被合并的 PR（文档/测试/feature）（演示）
\end{itemize}

\end{document}
